\documentclass[11pt,a4paper]{article}
\usepackage[T1]{fontenc}
\usepackage{lmodern}
\usepackage[margin=22mm]{geometry}
\usepackage{microtype}
\usepackage{graphicx}
\usepackage{booktabs}
\usepackage{xcolor}
\usepackage{listings}
\usepackage{enumitem}
\usepackage[hidelinks]{hyperref}
\setlength{\parindent}{0pt}
\setlength{\parskip}{6pt}
\setlist[itemize]{topsep=2pt, itemsep=2pt, parsep=0pt}
\lstset{language=Python, basicstyle=\ttfamily\footnotesize, breaklines=true,
        commentstyle=\color{gray}, stringstyle=\color{teal!70!black}, columns=fullflexible, showstringspaces=false}

% A figure that always stays exactly where it is written.
\newcommand{\plot}[2][\textwidth]{\begin{center}\includegraphics[width=#1]{figures/#2.png}\end{center}}

\begin{document}

{\Large\bfseries Assignment 2: Time series graphics}\par
{\small Data Forecasting. Plots made with Python (pandas + matplotlib); the full code is in the appendix and in \texttt{analysis.py}.}

\section*{Before we start}

Replacing a dataset to have seasonality.
None of the three Assignment 1 series is seasonal. NVIDIA and S\&P 500/VIX are financial market prices,
and the euro-area GDP was seasonally adjusted by Eurostat (code SCA), which means the seasonality had been removed on purpose.
As the assignment asks, dataset 3 is replaced by the same GDP series, not seasonally adjusted (code NSA):
same source (Eurostat table \texttt{namq\_10\_gdp}, real GDP index 2020 = 100, euro area EA20), same 120 quarters (1996 Q3 to 2026 Q2), no missing values.
It was downloaded on 21 October 2026 with \texttt{download\_gdp.py} and saved as \texttt{03\_euro\_area\_real\_gdp\_nsa\_quarterly.csv}.
This series is used from now on.

Number of lags. All lag plots use 12 lags (at least 9 were required): 12 lags cover one year of monthly data and three years of quarterly data,
so a yearly seasonal pattern would show up. For the weekly NVIDIA data, 12 lags cover about 3 months (a yearly pattern would need lag 52),
so its seasonality is judged from the line plot.

Words used:
\begin{itemize}
  \item Trend = a long-term increase or decrease.
  \item Seasonality = a pattern that repeats at a fixed period (e.g.\ every year).
  \item Cycle = rises and falls that are not of fixed length (e.g.\ recessions).
  \item Stationary = the series looks the same no matter when you look at it: no trend, no seasonality, roughly constant spread.
\end{itemize}

% =====================================================================
\newpage
\section{NVIDIA stock (weekly, 312 weeks, Oct 2020 to Oct 2026)}

Two variables: adjusted close price (USD) and weekly trading volume (shown in billions of shares).

\subsection*{Line plot}
\plot{1a_nvda_line}
Price:
\begin{itemize}
  \item Trend: strong upward trend, from about \$13 in 2020 to about \$234 in 2026, almost all of it after the start of 2023.
        Before that there is a rise in 2021 and a fall in 2022 (from about \$33 to \$11).
  \item Seasonality: none; nothing repeats at the same time each year.
  \item Cycle: no clear cycle. One full boom and bust (Nov 2021 to Oct 2022, $-66\%$); the later falls ($-22\%$ in summer 2024,
        $-36\%$ from Nov 2024 to Apr 2025) were short dips inside the uptrend.
  \item Variance: the swings get larger as the price gets higher (e.g.\ the fall from about \$147 in November 2024 to \$94 in April 2025), so the variance is not constant.
  \item Stationary: no (trend and growing variance).
\end{itemize}

Volume:
\begin{itemize}
  \item Trend: no steady trend, but the level changes: higher in 2022--2023 (around 2--3 billion shares a week),
        then a clear decline since mid-2024 down to about 0.5--1 billion.
  \item Seasonality: none visible.
  \item Cycle: no clear cycle.
  \item Variance: very noisy from week to week, with single-week spikes.
  \item Stationary: no, because the average level changes over time.
\end{itemize}

\subsection*{Seasonal plots}
Not made: the line plot shows no seasonality in either variable.

\subsection*{Scatter plot (price vs.\ volume)}
\plot[0.8\textwidth]{1b_nvda_scatter}
\begin{itemize}
  \item With two variables there is only one pair. The colour shows the year (light = 2020, dark = 2026).
  \item There is a negative relationship (correlation $-0.68$): high-price weeks have low volume.
  \item The colours show that this is mostly about time: in 2020--2023 (light dots) the price was low and the volume varied a lot,
        while in 2025--2026 (dark dots) the price is high and the volume is low.
  \item Both variables trend over time in opposite directions, so the plot does not prove that one causes the other.
  \item One likely reason: volume is counted in shares, and at a higher price fewer shares are needed for the same amount of money.
\end{itemize}

\subsection*{Lag plots: price}
\plot{1c_nvda_lag_price}
\begin{itemize}
  \item At every lag the dots lie very close to the diagonal (r goes only from 1.00 at lag 1 to 0.97 at lag 12):
        the price this week is almost the same as the price a few weeks ago.
  \item This strong, slowly fading dependence is typical of a trending, non-stationary series.
  \item At larger lags the cloud spreads out and moves above the diagonal, because the price now is usually higher than k weeks ago (the upward trend).
  \item No lag stands out, so there is no repeating pattern within 12 weeks.
\end{itemize}

\subsection*{Lag plots: volume}
\plot{1d_nvda_lag_volume}
\begin{itemize}
  \item A clear but much weaker positive relationship: r = 0.77 at lag 1, falling slowly to about 0.6 at lag 12.
  \item The wide cloud shows the large week-to-week noise.
  \item The slow decline in r comes from the long periods of high volume (2022--2023) and low volume (2025--2026):
        weeks inside the same period resemble each other.
  \item No lag stands out, so there is no repeating pattern within 12 weeks.
\end{itemize}

% =====================================================================
\newpage
\section{S\&P 500 and VIX (monthly, 108 months, Oct 2017 to Sep 2026)}

Three variables: S\&P 500 monthly mean close, VIX monthly mean close, VIX monthly maximum daily close.
VIX measures expected stock-market volatility (``fear index'').

\subsection*{Line plot}
\plot{2a_market_line}
S\&P 500:
\begin{itemize}
  \item Trend: clear upward trend, from about 2{,}600 to about 7{,}700 points.
  \item Seasonality: none; the falls listed below do not happen in the same month each year.
  \item Cycle: no regular cycle. The trend is interrupted by falls at irregular times: late 2018, the COVID crash in March 2020,
        a long decline during 2022, a drop in March--April 2025 and a smaller one in March 2026.
  \item Stationary: no (trend).
\end{itemize}

VIX (mean and max):
\begin{itemize}
  \item Trend: none; the series keeps returning to about 15--20.
  \item Seasonality: none; the spikes happen in different months.
  \item Cycle: no regular cycle. Sudden spikes fade within a few months: early 2018, end of 2018, March 2020 (the largest: mean 57.7, max 82.7), August 2024 and April 2025.
        From mid-2020 to the end of 2022 VIX stays above its usual level (mostly 20--30) without big spikes.
  \item Mean vs.\ max: the maximum is always above the mean (by definition), and the gap is largest during spikes.
  \item Stationary: roughly yes, the closest to stationary of all the series: no trend, no seasonality, a stable usual level; only the spikes make the spread uneven.
\end{itemize}

\subsection*{Seasonal plots}
Not made: the line plot shows no seasonality in any of the three variables.

\subsection*{Pairwise scatter plots}
\plot{2b_market_scatter}
\begin{itemize}
  \item All three pairs of variables; colour shows the year (light = 2017, dark = 2026).
  \item VIX mean vs.\ S\&P 500 and VIX max vs.\ S\&P 500 (left, middle): only a weak negative relationship (correlation $-0.14$ for both).
  \item The S\&P 500 trends up over time while VIX does not, so the S\&P level alone says little about VIX.
  \item But the highest VIX values occur in months when the S\&P 500 had just fallen (March 2020 at about 2{,}650 points, April 2025 at about 5{,}370).
  \item VIX max vs.\ VIX mean (right): a strong positive, almost straight-line relationship (correlation 0.95), as expected,
        since both are computed from the same daily values.
  \item The single point at the top right is March 2020.
\end{itemize}

\subsection*{Lag plots: S\&P 500}
\plot{2c_sp500_lag}
\begin{itemize}
  \item Very similar to the NVIDIA price: dots close to the diagonal at all lags (r from 0.99 down to 0.92), a sign of a trending, non-stationary series.
  \item As the lag grows, most dots move above the diagonal (the index is higher than k months ago, i.e.\ the upward trend).
  \item The group of dots below the diagonal around 4{,}000--4{,}700 is 2022, when the index was lower than in the months before.
  \item No lag stands out (lag 12 is not special), so no seasonality.
\end{itemize}

\subsection*{Lag plots: VIX mean}
\plot{2d_vix_mean_lag}
\begin{itemize}
  \item Very different from the S\&P 500. At lag 1 there is a clear positive relationship (r = 0.67): a high-volatility month is usually followed by another fairly high one.
  \item But it fades fast: r = 0.42 at lag 2, about 0.2 at lags 4--7 and about 0 at lags 9--12.
  \item So VIX only ``remembers'' the last few months, which fits a stationary series that returns to its usual level after each spike.
  \item Lag 12 shows nothing special (no yearly pattern), so no seasonality.
  \item The isolated dots far from the cloud come from March--April 2020.
\end{itemize}

\subsection*{Lag plots: VIX max}
\plot{2e_vix_max_lag}
\begin{itemize}
  \item Same picture as the VIX mean, slightly weaker: r = 0.61 at lag 1, about 0.1--0.2 at lags 3--8, near 0 at lags 9--12.
  \item The monthly maximum depends on a single day, so it is noisier than the monthly mean.
  \item Again a short memory, no seasonality.
  \item The dot at 82.7 is March 2020.
\end{itemize}

% =====================================================================
\newpage
\section{Euro-area real GDP, not seasonally adjusted (quarterly, 120 quarters, 1996 Q3 to 2026 Q2)}

One variable: real GDP index (2020 = 100). This is the seasonal replacement described at the start.

\subsection*{Line plot}
\plot{3a_gdp_line}
\begin{itemize}
  \item Trend: GDP grows from about 74 to about 115 over 30 years; growth is faster before 2008 than after.
  \item Seasonality: a saw-tooth that repeats every year (every 4 quarters): low in Q1, high in Q4.
  \item Cycle: the business cycle shows as irregular downturns: the 2008--2009 financial crisis, flat or falling GDP around 2011--2013 (euro debt crisis)
        and the COVID-19 collapse in 2020 Q2 followed by a fast rebound. These downturns are not at fixed intervals, so they are a cycle, not seasonality.
  \item Stationary: no (trend and seasonality).
\end{itemize}

\subsection*{Seasonal plot}
\plot[0.85\textwidth]{3b_gdp_seasonal}
\begin{itemize}
  \item Each line is one year (light = 1996, dark = 2026), plotted across its four quarters.
  \item Almost every line has the same shape: lowest in Q1, up in Q2, flat or slightly down in Q3, highest in Q4.
  \item In the 29 complete years (1997--2025), Q1 is the lowest quarter in 28 (all except 2020) and Q4 is the highest in 28 (all except 2008, the financial crisis).
  \item On average Q1 is about 2.4\% below its year's average and Q4 about 3.0\% above.
  \item The lines are stacked upwards because of the trend; the many lines bunched together around 90--98 are the slow-growth years 2008--2015.
  \item The one line that breaks the pattern is 2020, with its deep fall in Q2 (COVID-19 lockdowns).
\end{itemize}

\subsection*{Seasonal subseries plot}
\plot{3c_gdp_subseries}
\begin{itemize}
  \item Each panel shows one quarter across all years; the black line is that quarter's average.
  \item Q4 has the highest average (97.8) and Q1 the lowest (94.0), confirming the seasonal pattern.
        (The averages also contain the trend, so they are only a rough measure of the seasonal effect.)
  \item All four panels have the same shape (same trend, same dips in 2009 and 2020), which means the seasonal pattern is stable over time:
        it does not grow, shrink or change shape.
  \item The COVID-19 fall is by far the largest in Q2 (2020 Q2).
\end{itemize}

\subsection*{Lag plots}
\plot{3d_gdp_lag}
\begin{itemize}
  \item Dots are coloured by the quarter of y(t). All lags have a high r (0.89 to 0.97) because of the trend.
  \item At lag 1 the colours separate: Q1 dots (blue) lie below the diagonal (Q1 is lower than the Q4 just before it), and Q4 dots (green) lie above (Q4 is higher than Q3).
  \item At lags 4, 8 and 12 (the same quarter one, two and three years earlier) all colours mix in a narrow band along the diagonal,
        and r is higher than at the neighbouring lags (lag 4: 0.97 vs.\ 0.94 at lag 3 and 0.92 at lag 5; lag 8: 0.94 vs.\ 0.91 at lag 7 and 0.90 at lag 9).
  \item Together, this is the signature of seasonality with period 4 quarters (one year).
  \item In every panel, two dots lie far from the diagonal, and both come from the single COVID-19 quarter 2020 Q2 (GDP 91.4).
        The dot below the line is 2020 Q2 itself as the value now (always orange).
        The dot above the line is the quarter k periods later, compared with that very low value; its colour changes with the lag
        (e.g.\ lag 1: 2020 Q3, pink; lag 4: 2021 Q2, orange).
\end{itemize}

% =====================================================================
\newpage
\section{Summary}

\begin{center}\small
\begin{tabular}{@{}lcccc@{}}
\toprule
Series & Trend & Seasonality & Cycle & Looks stationary? \\
\midrule
NVIDIA price             & yes, up        & no  & no (one boom/bust, short dips)  & no \\
NVIDIA volume            & level changes  & no  & no                  & no \\
S\&P 500                 & yes, up        & no  & irregular falls     & no \\
VIX mean / VIX max       & no             & no  & irregular spikes    & roughly yes \\
Euro-area GDP (NSA)      & yes, up        & yes, yearly & yes (recessions) & no \\
\bottomrule
\end{tabular}
\end{center}

\begin{itemize}
  \item Most of the series are not stationary because of a trend; to model them, differencing (or returns) would be needed. VIX is the only one that looks roughly stationary.
  \item The euro-area GDP (not seasonally adjusted) is the seasonal series that will be used from now on: a stable yearly pattern with Q1 low and Q4 high.
  \item In the multivariate datasets, the strongest relationship is VIX max vs.\ VIX mean. The other relationships (NVIDIA price vs.\ volume, S\&P 500 vs.\ VIX) are weaker and mostly driven by time.
\end{itemize}

% =====================================================================
\newpage
\appendix
\section*{Appendix: code}

How to run. From the assignment folder: \texttt{pip install -r requirements.txt}, then \texttt{python analysis.py}.
It reads the three CSV files and saves all the plots in \texttt{figures/}.
\texttt{download\_gdp.py} was run once to create the new seasonal GDP CSV.

\subsection*{\texttt{analysis.py}}
\lstinputlisting{analysis.py}

\subsection*{\texttt{download\_gdp.py}}
\lstinputlisting{download_gdp.py}

\end{document}
